\chapter{Estado da arte}\label{cap:estado-da-arte}

% Manual: 4 a 5 trabalhos dos últimos 5 a 10 anos. Para cada um, responder:
% qual problema tentou resolver, como abordou (algoritmos, técnicas, arquitetura),
% quais foram os principais achados e o que faltou. Depois, comparar os trabalhos
% entre si e com este TCC, fechar com o quadro comparativo e deixar clara a lacuna.
% Ler cada artigo antes de escrever e conferir os números no PDF.

\pendente{Parágrafo de abertura: critérios de busca (bases, período, palavras-chave) e de seleção dos trabalhos.}

\section{Decisão de colheita do café com rede neural convolucional}

% mundimfilho2021 (artigo do SBSI) e santos2022 (dissertação da mesma linha, UFU).
% Antecedente mais próximo: aplicativo Android offline, em português, que responde
% "colher" ou "não colher". Deixar explícito em que este TCC se diferencia.
\pendente{A escrever.}

\section{Detecção de frutos de café em ramos com YOLO}

% bazame2023: fotos de smartphone em Patos de Minas-MG; comparação entre
% versões do YOLO e efeito do tamanho da imagem de entrada.
\pendente{A escrever.}

\section{Aplicativo de maturação de azeitonas com YOLOv8}

% alvino2024: aplicativo brasileiro (Flutter, API em Flask e YOLOv8) que depende de servidor.
\pendente{A escrever.}

\section{Detecção leve da maturação do tomate}

% dong2025: GPC-YOLO, tomate em ambiente não estruturado, 3 classes; não foi testado em celular.
\pendente{A escrever.}

\section{Aplicativo Android com YOLO leve executado no dispositivo}

% zheng2025: arroz; modelo embarcado com testes de desempenho no próprio smartphone.
% Bom modelo de metodologia para a avaliação no celular.
\pendente{A escrever.}

\section{Quadro comparativo e posicionamento deste trabalho}

\pendente{Parágrafo comparando as abordagens (diálogo entre os autores) e deixando clara a lacuna que este trabalho preenche.}

\begin{table}[htb]
  \ABNTEXfontereduzida
  \caption{Comparação entre os trabalhos correlatos e este trabalho}
  \label{tab:comparativo}
  \begin{tabularx}{\textwidth}{lXXXX}
    \hline
    \textbf{Trabalho} & \textbf{Cultura} & \textbf{Modelo} & \textbf{Onde roda} & \textbf{Saída para o produtor} \\
    \hline
    \citeonline{mundimfilho2021} & Café & CNN (MobileNet) & \pendente{conferir} & Colher ou não colher \\
    \citeonline{bazame2023} & Café & YOLOv3 e YOLOv4 & \pendente{conferir} & \pendente{conferir} \\
    \citeonline{alvino2024} & Azeitona & YOLOv8 & Servidor & Índice de maturação \\
    \citeonline{dong2025} & Tomate & GPC-YOLO (YOLOv8n modificado) & \pendente{conferir} & \pendente{conferir} \\
    \citeonline{zheng2025} & Arroz & YOLO leve & No smartphone & \pendente{conferir} \\
    Este trabalho & Tomate & YOLO26n-cls & No celular, offline & Colher ou esperar \\
    \hline
  \end{tabularx}
  \legend{Fonte: elaborada pelos autores.}
\end{table}
